\section{Mandelbrot recurrence in real arithmetic}
\label{sec:real-form}

CUDA arithmetic operates on real scalar components, so the complex recurrence must
first be expressed as additions and multiplications of real numbers. Expanding the
square gives
\[
  (x+iy)^2=x^2+2ixy+i^2y^2=(x^2-y^2)+i(2xy),
\]
because $i^2=-1$. Matching real and imaginary components produces the update below.

Writing \(z = x + i y\) and \(c = x_0 + i y_0\), the iteration becomes:

\begin{align}
x_{n+1} &= x_n^2 - y_n^2 + x_0, \\
y_{n+1} &= 2 x_n y_n + y_0.
\end{align}

Both updates consume $x_n,y_n$ from the same iteration. The new real value must be
kept in a temporary until the imaginary update has finished using $x_n$.

\paragraph{Worked complex update.}
For $z_n=1+i$ and $c=-1/2+i/4$, the square is $2i$, so
$z_{n+1}=-1/2+9i/4$. The real formulas give the same result:
$1^2-1^2-1/2=-1/2$ and $2\cdot1\cdot1+1/4=9/4$.
Its squared magnitude is $1/4+81/16=85/16>4$.

The escape test is:

\begin{equation}
x_n^2 + y_n^2 > 4 \quad\Rightarrow\quad \text{escaped.}
\end{equation}

The strict inequality is the mathematical escape guarantee. Equality alone is
insufficient: $c=-2$ generates $0,-2,2,2,\ldots$, a bounded orbit. Rendering kernels
can use a different comparison or a larger bailout for numerical and counting
conventions; those choices must be distinguished from the theorem. Reaching an
iteration limit without escape does not in general prove membership in the set.

Many kernels cache squares:
\[
zrsqr = x^2,\quad zisqr = y^2,
\]
so that \(x^2-y^2\) and \(x^2+y^2\) can be formed cheaply as \code{zrsqr - zisqr}
and \code{zrsqr + zisqr}. This optimization is especially valuable when \(x\) and \(y\) are
represented by multi-component expansion types.
